% !TEX root = all_beamer.tex
%

\subtitle{Course of Machine Learning}
\author{Giorgio Gambosi}
\institute[Tor Vergata]{Master's Degree in Computer Science\\University of Rome ``Tor Vergata''}
\date{A.Y. 2025--2026}

\usepackage{tikz}
\usepackage{tikz-network}
\usepackage[customcolors]{hf-tikz}
\usepackage{pgfplots}
\pgfplotsset{compat=1.18}
\usetikzlibrary{arrows, decorations.pathmorphing, positioning, backgrounds, fit, shadows, automata, patterns}
\usetikzlibrary{calc, intersections, through, shapes, scopes, chains, matrix, shapes.arrows, decorations, snakes}
% ---------- Theme & colours ----------
\usetheme{Madrid}
\usecolortheme{whale}
%\useoutertheme{miniframes}
\useinnertheme{rounded}
\setbeamertemplate{navigation symbols}{}

\definecolor{MainBlue}{RGB}{31,56,100}
\definecolor{AccentBlue}{RGB}{70,130,180}
\definecolor{AlertRed}{RGB}{160,30,30}
\definecolor{ExGreen}{RGB}{30,120,60}
\definecolor{BlockBg}{RGB}{220,232,246}
\definecolor{torgray}{RGB}{220,225,235}

\setbeamercolor{frametitle}{bg=MainBlue, fg=white}
\setbeamercolor{title}{bg=MainBlue, fg=white}
\setbeamercolor{block title}{bg=AccentBlue, fg=white}
\setbeamercolor{block body}{bg=BlockBg}
\setbeamercolor{block title alerted}{bg=AlertRed, fg=white}
\setbeamercolor{block body alerted}{bg=red!8}
\setbeamercolor{block title example}{bg=ExGreen, fg=white}
\setbeamercolor{block body example}{bg=green!7}
\setbeamercolor{structure}{fg=MainBlue}

% ---------- Fonts ----------
\usefonttheme{professionalfonts}
\setbeamerfont{frametitle}{size=\normalsize,series=\bfseries}
\setbeamerfont{block title}{size=\small,series=\bfseries}

% ---------- Packages ----------
\usepackage{amsmath, amssymb, mathtools, amsthm}
\usepackage{booktabs}
\usepackage{tikz}
\usetikzlibrary{arrows.meta, positioning, shapes.geometric, fit}
\usepackage{xcolor}
\usepackage{microtype}
\usepackage{graphicx}

% ---------- Theorem styles ----------
\setbeamertemplate{theorems}[numbered]
\newtheorem{mytheorem}{Theorem}
\newtheorem{mydefinition}{Definition}

% ---------- Macros ----------

\providecommand{\calX}{\mathcal{X}}
\providecommand{\calY}{\mathcal{Y}}
\providecommand{\calH}{\mathcal{H}}
\providecommand{\calT}{\mathcal{T}}
\providecommand{\calA}{\mathcal{A}}
\providecommand{\calR}{\mathcal{R}}
\providecommand{\calF}{\mathcal{F}}
\providecommand{\barR}{\overline{\mathcal{R}}}
\providecommand{\pM}{p_M}
\providecommand{\pC}{p_C}
\providecommand{\indicator}[1]{\mathbf{1}\!\left[#1\right]}
\providecommand{\argmax}{\operatorname*{arg\,max}}
\providecommand{\argmin}{\operatorname*{arg\,min}}
\providecommand{\vcd}[1]{\operatorname{VCdim}(#1)}
\providecommand{\Prob}[2]{\mathbb{P}_{#1}\!\left[#2\right]}

\usepackage{xkcdcolors}
%
%
%\definecolor{primary}  {RGB}{31,  56, 100}
%\definecolor{accent}   {RGB}{46, 117, 182}
%\definecolor{lightbg}  {RGB}{235,242,252}
%\definecolor{alertred} {RGB}{180, 30,  30}
%\definecolor{greentick}{RGB}{0,  128,  64}
%
%%% node colours matching the original xkcd names
\definecolor{nodeblue}  {RGB}{101,140,196}  % xkcdFadedBlue
%\definecolor{nodered}   {RGB}{198, 42, 86}  % xkcdLipstickRed
%\definecolor{nodeturq}  {RGB}{  6,194,172}  % xkcdTurquoise
%\definecolor{nodeorange}{RGB}{249,149, 40}  % xkcdOrange
%\definecolor{nodeameth} {RGB}{155, 95,192}  % xkcdAmethyst
%\definecolor{nodeteal}  {RGB}{  2,130,122}  % xkcdTeal
%\definecolor{nodegrey}  {RGB}{190,190,190}  % xkcdLightGrey
%
%\setbeamercolor{palette primary}   {bg=primary,  fg=white}
%\setbeamercolor{palette secondary} {bg=accent,   fg=white}
%\setbeamercolor{palette tertiary}  {bg=primary,  fg=white}
%\setbeamercolor{palette quaternary}{bg=primary,  fg=white}
%\setbeamercolor{structure}         {fg=primary}
%\setbeamercolor{section in toc}    {fg=primary}
%\setbeamercolor{frametitle}        {bg=primary,  fg=white}
%\setbeamercolor{title}             {fg=white}
%\setbeamercolor{block title}       {bg=accent,   fg=white}
%\setbeamercolor{block body}        {bg=lightbg}
%\setbeamercolor{block title alerted}{bg=alertred, fg=white}
%\setbeamercolor{block body alerted} {bg=lightbg}
%\setbeamercolor{block title example}{bg=greentick,fg=white}
%\setbeamercolor{block body example} {bg=lightbg}
%\setbeamercolor{alerted text}       {fg=alertred}
%\setbeamercolor{itemize item}       {fg=accent}
%\setbeamercolor{itemize subitem}    {fg=primary}
%
%\setbeamertemplate{itemize item}   {\small\raise1pt\hbox{$\bullet$}}
%\setbeamertemplate{itemize subitem}{\small\raise1pt\hbox{$-$}}
%\setbeamertemplate{navigation symbols}{}
%\setbeamertemplate{section in toc}[sections numbered]
%
%%% ── Packages ────────────────────────────────────────────────────────────────
%\usepackage[utf8]{inputenc}
%\usepackage[T1]{fontenc}
%\usepackage{amsmath,amssymb}
%\usepackage{array}
%\usepackage{tikz}
%\usetikzlibrary{arrows.meta,positioning,calc}
%
%%% ── TikZ diagram styles ─────────────────────────────────────────────────────
%\tikzset{
%  dnode/.style={
%    rectangle, draw=black!25, line width=0.4pt,
%    fill=#1, fill opacity=0.50, text opacity=1,
%    font=\small, inner sep=4pt, minimum height=7mm
%  },
%  darrow/.style={
%    -{Stealth[length=4pt,width=3pt]},
%    line width=0.5pt, draw=black!45
%  },
%  dlabel/.style={font=\footnotesize\itshape, text=black!55}
%}
%
%%% ── Fonts ───────────────────────────────────────────────────────────────────
%\setbeamerfont{frametitle}  {size=\large,    series=\bfseries}
%\setbeamerfont{title}       {size=\Large,    series=\bfseries}
%\setbeamerfont{block title} {size=\normalsize,series=\bfseries}
%
%%% ── Custom math / text commands ─────────────────────────────────────────────
\providecommand{\term}[1]{\textbf{\color{accent}#1}}
\providecommand{\Xmat}{\mathbf{X}}
%\providecommand{\xvec}{\mathbf{x}}
\providecommand{\wvec}{\mathbf{w}}
\providecommand{\ind}[1]{\mathbf{1}[#1]}
\providecommand{\Rr}{\mathcal{R}}
\providecommand{\Tr}{\mathcal{T}}
\providecommand{\Hr}{\mathcal{H}} 

%\usepackage[english]{babel}
%\RequirePackage{etex}
%
%\usepackage[utf8]{inputenc}
%\usepackage[T1]{fontenc}
%\usepackage{CormorantGaramond}
%\usepackage{geometry}
%\usepackage{xcolor}
%\usepackage{xkcdcolors}
%\usepackage{graphicx}
%\usepackage{csquotes}
%\usepackage{xstring}
%
%% --- MATEMATICA AVANZATA E SIMBOLI ---
%\usepackage{amsmath, amsfonts, amssymb, amsthm, mathtools}
%\usepackage{dsfont, mathrsfs, relsize, mleftright}
\usepackage{bm}            % Bold math
\usepackage{upgreek}       % Greche dritte
\usepackage{derivative}    % Gestione derivate semplificata
%\usepackage{siunitx}       % Unità di misura SI
%\usepackage{textgreek}     % Simboli greci nel testo
%\usepackage[misc]{ifsym}   % Per \Letter e simboli vari
%\usepackage{cancel}        % Per barrare termini nelle equazioni
%\usepackage{centernot}     % Per negazioni centrate
%\usepackage{annotate-equations}
%
%% --- TABELLE E LAYOUT ---
%\usepackage{booktabs, tabularx, longtable, multirow, colortbl}
%\usepackage{pdflscape, rotating}
%\usepackage{multicol}
%\usepackage{enumitem}
%\usepackage{subcaption}
%\usepackage{environ}
%\usepackage[skins, breakable]{tcolorbox}
%\usepackage{snotez}        % Per sidenote
%
%% --- ALGORITMI E CODICE ---
%\usepackage[algoruled]{algorithm2e}
%\usepackage{algorithmic}
%\usepackage{listings}
%
%% --- TIKZ E PGFPLOTS ---
\usepackage{tikz}
\usepackage{tikz-network}
\usepackage[customcolors]{hf-tikz}
\usepackage{pgfplots}
\pgfplotsset{compat=1.18}
\usetikzlibrary{arrows, decorations.pathmorphing, positioning, backgrounds, fit, shadows, automata, patterns}
\usetikzlibrary{calc, intersections, through, shapes, scopes, chains, matrix, shapes.arrows, decorations, snakes}

\tikzset{
  >=latex, % for default LaTeX arrow head
  node/.style={thick,circle,draw=myblue,minimum size=22,inner sep=0.5,outer sep=0.6},
  node in/.style={node,green!20!black,draw=mygreen!30!black,fill=mygreen!25},
  node hidden/.style={node,blue!20!black,draw=myblue!30!black,fill=myblue!20},
  node convol/.style={node,orange!20!black,draw=myorange!30!black,fill=myorange!20},
  node out/.style={node,red!20!black,draw=myred!30!black,fill=myred!20},
  connect/.style={thick,mydarkblue}, %,line cap=round
  connect arrow/.style={-{Latex[length=4,width=3.5]},thick,mydarkblue,shorten <=0.5,shorten >=1},
  node 1/.style={node in}, % node styles, numbered for easy mapping with \nstyle
  node 2/.style={node hidden},
  node 3/.style={node out}
}
%
%% --- CONFIGURAZIONE PAGINA E GEOMETRY ---
%\geometry{
%    left=2.5cm,
%    right=2.5cm,
%    top=3cm,
%    bottom=3.5cm,
%    marginparwidth=0pt,
%    marginparsep=0pt
%}
%
%%\makeatletter
%%%\renewcommand{\chaptermark}[1]{\markboth{#1}{}}
%%
%%\def\ppfooterB{\footnotesize  \@leadauthor\hspace{7pt}|\hspace{7pt}\@title\hspace{7pt}|\hspace{7pt}\thepage\space of\space\pageref{LastPage}}
%%%\def\ppfooterC{\footnotesize \@leadauthor\hspace{7pt}|\hspace{7pt}\@shorttitle}
%%\def\ppfooterC{{\footnotesize 
%%    \ifnum\value{chapter}>0 Cap. \thechapter\hspace{7pt}|\hspace{7pt}\leftmark\fi 
%%}}
%%
%%\lfoot{\ppfooterC}
%%\rfoot{\ppfooterB}
%%\preto{\footrule}{\color{MediumGrey}}
%%
%%\fancypagestyle{plain}{
%%    \fancyhf{}
%%    \fancyfoot[R]{\ppfooterB}
%%    \fancyfoot[L]{\ppfooterC}
%%}
%%\makeatother
%
%% --- BIBLIOGRAFIA ---
%\usepackage[
%    backend=biber,
%    style=authoryear,
%    natbib=true,
%    hyperref=true,
%    alldates=year
%]{biblatex}
%
%% --- DEFINIZIONE COLORI (Senza duplicati) ---
%\definecolor{node}{RGB}{220,220,220}
%\definecolor{box}{RGB}{240,240,240}
%\definecolor{ft}{RGB}{210,40,40}
%\definecolor{bt}{RGB}{54,94,155}
%\definecolor{nt}{RGB}{90,90,90}
%\definecolor{mt}{RGB}{168,27,49}
%\definecolor{at}{RGB}{170,80,80}
%
%% Colori specifici Berkeley/Metropolis
%\definecolor{BerkeleyGold}{RGB}{193,129,45}
%\definecolor{BerkeleyBlueText}{RGB}{48,99,126}
%\definecolor{mDarkBrown}{HTML}{604c38}
%\definecolor{mDarkTeal}{HTML}{23373b}
%
%% Configurazione semantica Colorlet
%\colorlet{pale}{xkcdLightGrey}
%\colorlet{observed}{xkcdPinkish!60}
%\colorlet{param}{xkcdLightTan}
%\colorlet{latent}{xkcdPeriwinkle}
%\colorlet{evidmath}{xkcdBlack}
%\colorlet{odata}{xkcdDuskBlue}
%\colorlet{elat}{xkcdPurplishRed}
%\colorlet{eparam}{xkcdGreenishBlue}
%\colorlet{edummy}{xkcdBurntSiena}
%\colorlet{efunc}{xkcdBurntSiena}
%\colorlet{evar}{xkcdLightishRed}
%\colorlet{evidmath}{xkcdBlack}
%\colorlet{lightmath}{xkcdMediumGrey!80}
%\colorlet{odata}{xkcdPetrol!80}
%\colorlet{emath1}{xkcdBrightRed}
%\colorlet{emath2}{xkcdRichPurple}
%\colorlet{emath}{xkcdGreenishBlue}
%\colorlet{edummy}{xkcdBurntSiena}
%\colorlet{efunc}{xkcdBurntSiena}
%\colorlet{efunc1}{xkcdNeonPink}
%\colorlet{evar}{xkcdRuby}
%\colorlet{eparam}{xkcdDuskBlue}
%\colorlet{elat}{xkcdPurplishRed}
%\colorlet{egen}{xkcdAmethyst}
%
\colorlet{myred}{red!80!black}
\colorlet{myblue}{blue!80!black}
\colorlet{mygreen}{green!60!black}
\colorlet{myorange}{orange!70!red!60!black}
\colorlet{mydarkred}{red!30!black}
\colorlet{mydarkblue}{blue!40!black}
\colorlet{mydarkgreen}{green!30!black}
%
%% --- AMBIENTI PERSONALIZZATI ---
%\tcbset{enhanced jigsaw, colback=xkcdTeal!5!white, colframe=xkcdTeal!75!black, fonttitle=\bfseries}
%
%\newenvironment{notebox}[1]{
%    \medskip\noindent
%    \begin{tcolorbox}[breakable, width=\textwidth, boxrule=.1mm, title=#1]
%    \color{xkcdBattleshipGrey!75!black}
%}{\end{tcolorbox}}
%
%\newenvironment{myequation}{\begin{equation*}\color{evidmath}}{\end{equation*}}
%\NewEnviron{myalign}{{\color{evidmath}\begin{align*}\BODY\end{align*}}}
%\newenvironment{myexample}{\bigskip\hrule\medskip\small}{\medskip\hrule\bigskip}
%
%\newtheorem{exmp}{Example}[section]
%%\newtheorem{definition}{Definition}
%%\renewtheorem{theorem}{Theorem}
%
%% --- MACRO ALGEBRA LINEARE (Vettori e Matrici) ---
\providecommand{\Vbold}[1]{\ifmmode\mathbf{#1}\else\textbf{#1}\fi}
\providecommand{\Zero}{\Vbold{0}}
\providecommand{\Vuno}{\Vbold{1}}

% Generazione automatica lettere Maiuscole Bold
\makeatletter
% Lista delle maiuscole
\@for\letter:=A,B,C,D,F,H,I,K,N,P,R,S,T,U,V,W,X,Y,Z\do{%
  \expandafter\edef\csname\letter\endcsname{\noexpand\Vbold{\letter}}%
}

% Lista delle minuscole
\@for\lcase:=a,b,l,n,u,y,w,z,x,t,v\do{%
  \expandafter\edef\csname\lcase\endcsname{\noexpand\Vbold{\lcase}}%
}
\makeatother
\providecommand{\TT}{\Vbold{T}}

%% Lettere Minuscole Bold
%
%
%% Greche Bold (upgreek/amsmath)
\providecommand{\VGamma}{\boldsymbol{\Upgamma}}
\providecommand{\Vlambda}{\boldsymbol{\lambda}}
\providecommand{\VLambda}{\boldsymbol{\Uplambda}}
\providecommand{\VSigma}{\boldsymbol{\Upsigma}}
\providecommand{\Vmu}{\boldsymbol{\mu}}
\providecommand{\Vtheta}{\boldsymbol{\theta}}
\providecommand{\VTheta}{\boldsymbol{\Theta}}
\providecommand{\Vphi}{\boldsymbol{\phi}}
\providecommand{\VPhi}{\boldsymbol{\Phi}}
\providecommand{\Vpsi}{\boldsymbol{\psi}}
\providecommand{\VPsi}{\boldsymbol{\Psi}}
\providecommand{\Vpi}{\boldsymbol{\pi}}
\providecommand{\Valpha}{\boldsymbol{\alpha}}
\providecommand{\Vbeta}{\boldsymbol{\beta}}
\providecommand{\Vxi}{\boldsymbol{\xi}}
\providecommand{\Vepsilon}{\boldsymbol{\epsilon}}
%
%\providecommand{\oa}{\overline{\a}}
%\providecommand{\ou}{\overline{\u}}
%\providecommand{\ov}{\overline{\v}}
%\providecommand{\ow}{\overline{\w}}
%\providecommand{\ox}{\overline{\x}}
%\providecommand{\oy}{\overline{\y}}
%\providecommand{\oz}{\overline{\z}}
%\providecommand{\oX}{\overline{\X}}
%\providecommand{\oVxi}{\overline{\Vxi}}
%\providecommand{\oVbeta}{\overline{\Vbeta}}
%\providecommand{\oVgamma}{\overline{\Vgamma}}
%\providecommand{\oVphi}{\overline{\Vphi}}
%\providecommand{\oY}{\overline{\Y}}
%\providecommand{\oBeta}{\overline{\B}}
%
%\providecommand{\func}[1]{{\color{efunc}#1}}
%\providecommand{\prm}[1]{#1}
%\providecommand{\lat}[1]{{\color{elat}\overline #1}}
%\providecommand{\val}[1]{#1}
%\providecommand{\vr}[1]{#1}
%
%\providecommand{\vlX}{\X}
%\providecommand{\vlxi}{\x_i}
%\providecommand{\vrt}{\Vtheta}
%\providecommand{\vrp}{\Psi}
%\providecommand{\lz}{{\color{elat}\overline\z}}
%\providecommand{\lzi}{\lat{\z_i}}
%
%
%%\providecommand{\ox}{{\color{evidmath}\overline\x}}
%\providecommand{\vlx}{\ox}
%\providecommand{\prts}{{\color{eparam}\Vtheta^*}}
%\providecommand{\prt}{{\color{eparam}\hat\Vtheta}}
%%\providecommand{\prm}[1]{{\color{eparam}#1}}
%%\providecommand{\X}{{\color{evidmath}\mathbf{X}}}
%\providecommand{\vlxu}{{\color{evidmath}\overline\x_1}}
%%\providecommand{\vlxi}{{\color{evidmath}\overline\x_i}}
%\providecommand{\vlxn}{{\color{evidmath}\overline\x_n}}
%%\providecommand{\Z}{{\color{elat}\mathbf{Z}}}
%\providecommand{\dummy}[1]{{\color{lightmath}#1}}
%\providecommand{\dz}{\dummy{\z}}
%
%
%\providecommand{\hy}{\hat{y}}
%\providecommand{\tm}{\tilde{\m}}
%\providecommand{\ty}{\tilde{y}}
%\providecommand{\tx}{\tilde{\x}}
%\providecommand{\tX}{\tilde{\X}}
%
%\providecommand{\vx}{\mathbf{x}}
%\providecommand{\vX}{\mathbf{X}}
%\providecommand{\vZ}{\mathbf{Z}}
%\providecommand{\vtheta}{\boldsymbol{\theta}}
%\providecommand{\vthetastar}{\boldsymbol{\theta}^*}
%\providecommand{\vxi}{\mathbf{x}_i}
%
%% Shortcut Caligrafici
%\providecommand{\cX}{\mathcal{X}}
%\providecommand{\cY}{\mathcal{Y}}
%\providecommand{\cH}{\mathcal{H}}
%\providecommand{\cXY}{\mathcal{X}\times\mathcal{Y}}
%\providecommand{\oR}{\overline{\mathcal{R}}}
%\providecommand{\cA}{\mathcal{A}}
%\providecommand{\vcd}[1]{\text{VCdim}(#1)}
%
%% --- STATISTICA E PROBABILITÀ ---
%\DeclareMathOperator{\E}{\mathbb{E}}
%\DeclareMathOperator*{\PP}{\mathds{P}}
%\DeclareMathOperator*{\EE}{\mathds{E}}
%
%\providecommand{\me}[1]{\E \left[ #1 \right]}
%\providecommand{\med}[2]{\E_{#1}\left[ #2 \right]}
%\providecommand{\ev}[2]{\EE_{#1}\left[ #2 \right]}
%\providecommand{\prob}[2]{\PP_{#1}\left[ #2 \right]}
%
%\providecommand{\var}[1]{\text{\textit{Var}}[#1]}
%\providecommand{\cov}[1]{\text{\textit{Cov}}[#1]}
%\providecommand{\diag}[1]{\text{\textit{diag}}#1}
%
%\providecommand{\gaussian}[3]{\frac{1}{\sqrt{2\pi}#3}e^{-\frac{(#1-#2)^{2}}{2#3^{2}}}}
%\providecommand{\gaussiandistr}[3]{#1\sim{\mathcal{N}}(#2,#3)}
%
%% --- OPERATORI E CALCOLO ---
%\providecommand{\bydef}{\overset{\Delta}{=}}
%\providecommand{\defeq}{\bydef}
\providecommand{\Real}{\mathbb{R}}
%\providecommand{\Integer}{\mathbb{I}}
%\providecommand{\Natural}{\mathbb{N}}
%\providecommand{\Complex}{\mathbb{C}}
%\providecommand{\vnorm}[1]{\left |\left | #1\right |\right |}
%
%\providecommand{\subfunc}[3]{\underset{#2}{\textrm{#1}}\:#3}
%\providecommand{\argmin}[2]{\subfunc{argmin}{#1}{#2}}
%\providecommand{\argmax}[2]{\subfunc{argmax}{#1}{#2}}
%
%\providecommand{\fracpartial}[2]{\pdv{#1}{#2}}
%\providecommand{\gradient}[2]{\mathlarger{\mathlarger{\nabla}}_{\mathsmaller{#1}} #2}
%
%\providecommand{\tset}{{\cal T}}
%
%\providecommand{\matr}[1]{\boldsymbol{#1}}
%
%%matrice righe
%\providecommand{\matrighe}[2]{
% \left(
%     \begin{array}{ccc}
%          \mbox{--} &  #1  &  \mbox{--} \\
%          & \vdots &  \\
%           \mbox{--} &  #2  &  \mbox{--} 
%     \end{array}
%\right)
%}
%     
%%matcolonne
%\providecommand{\matcolonne}[2]{
%     \left(
%     \begin{array}{ccc}
%          \vert  &  &  \vert \\
%         #1 & \cdots & #2 \\
%           \vert  &  &  \vert 
%     \end{array}
%     \right)}
%     
% \providecommand{\vcol}[2]{
%\left(
%     \begin{array}{c}
%          #1 \\
%          \vdots \\
%          #2
%     \end{array}
%     \right)}
%     
%\providecommand{\vcoll}[3]{
%\left(
%     \begin{array}{c}
%          #1 \\
%          #2 \\
%          \vdots \\
%          #3
%     \end{array}
%     \right)}
%
%\providecommand{\vriga}[2]{
%\left(
%\begin{array}{ccc}
%#1 & \cdots & #2
%\end{array}
%\right)
%}
%
%\providecommand{\vrigal}[3]{
%\left(
%\begin{array}{cccc}
%#1 & #2 & \cdots & #3
%\end{array}
%\right)
%}
%
%\providecommand{\cseq}[3]{#1,#2,\ldots ,#3}
%
%% somma 1,2,...,n
%\providecommand{\seq}[4]{#1#4 #2#4\ldots #4#3}
%
%% elenco x_1,x_2,...,x_n
%\providecommand{\iseq}[2]{#1_{1},#1_{2},\ldots ,#1_{#2}}
%
%% elenco x_1,...,x_n
%\providecommand{\iseqq}[2]{#1_{1},\ldots ,#1_{#2}}
%
%\providecommand{\tr}[1]{\text{tr}(#1)}
%% --- INDIPENDENZA STOCASTICA ---
%\providecommand{\CI}{\mathrel{\text{\scalebox{1.07}{$\perp\mkern-8mu\perp$}}}}
%\providecommand{\nCI}{\cancel{\CI}}
%\def\cind#1#2#3{#1\CI #2 \mid #3}
%\def\ind#1#2{#1\CI #2}
%
%% --- MACRO VAE / LATENTI (Colorate) ---
%%\providecommand{\val}[1]{{\color{odata}{#1}}}
%%\providecommand{\lat}[1]{{\color{elat}{#1}}}
%%\providecommand{\prm}[1]{{\color{eparam}{#1}}}
%
%%\providecommand{\vlx}{\val{\x}} 
%%\providecommand{\vlX}{\val{\X}}
%%\providecommand{\lz}{\lat{\z}} 
%%\providecommand{\lZ}{\lat{\Z}}
%%\providecommand{\prt}{\prm{\Vtheta}}
%%\providecommand{\prp}{\prm{\Vphi}}
%\DeclareMathOperator*{\HH}{\mathds{H}}
%\providecommand{\entr}[2]{\HH_{\scalebox{.8}{$\scriptscriptstyle #1$}}\mleft[\,#2\,\mright]}
%\providecommand{\dkl}[2]{D_{KL} \left({#1} || {#2} \right) }
%
%
%% --- CONFIGURAZIONE LISTINGS (Python) ---
%\lstset{
%    language=Python,
%    basicstyle=\ttfamily\color{xkcdDarkAquamarine},
%    keywordstyle=\ttfamily\color{xkcdRedOrange},
%    stringstyle=\color{green!80!black},
%    showstringspaces=false            
%}
%
%% --- TIKZ NEURAL NETWORKS ---
%\providecommand\denselyConnectNodes[2]{
%  \foreach\n [count=\lyrIdx, remember=\lyrIdx as \previdx] in #2 {
%    \foreach\y in {1,...,\n} {
%      \ifnum\lyrIdx>1
%        \foreach\x in {1,...,\prevn}
%          \draw[->] (#1-\previdx-\x) -- (#1-\lyrIdx-\y);
%      \fi
%    }
%    \xdef\prevn{\n}
%  }
%}
%
%% Setup Grafica finale
%\graphicspath{{graphics/}}
%\setkeys{Gin}{width=\linewidth,totalheight=\textheight,keepaspectratio}
%
%% Importazione comandi esterni (se ancora necessari)
%%\usepackage{commands} 
%\definecolor{titlecolor}{named}{xkcdDarkAquamarine}
%\definecolor{alertcolor}{named}{xkcdDuskyRose}
%
%\definecolor{evidmath}{named}{xkcdCadetBlue}
%\definecolor{lightmath}{named}{xkcdMediumGrey}
%%\definecolor{emath1}{named}{xkcdRedOrange}
%\definecolor{emath1}{named}{xkcdPumpkin}
%\definecolor{emath2}{named}{xkcdIrishGreen}
%\definecolor{emath3}{named}{xkcdDuskyRose}
%
%\providecommand{\gaussiansimplest}[1]{\frac{1}{\sqrt{2\pi}}\textit{exp}\left(-\frac{#1^{2}}{2}\right)}
%
%%%gaussiana multivariata
%\providecommand{\mgaussian}[4]{\frac{1}{(2\pi)^{#4/2}|#3|^{1/2}}\textit{exp}\left(-\frac{1}{2}(#1-#2)^{T}#3^{-1}(#1-#2)\right)}
%
%
%\providecommand{\mL}{\mathcal{L}}
%\providecommand{\mK}{\mathcal{K}}
%\providecommand{\mZ}{\mathcal{Z}}
%\providecommand{\mP}{\mathcal{P}}
%\providecommand{\mQ}{\mathcal{Q}}
%
%\renewcommand{\alert}[1]{{\bf\color{xkcdDuskyRose}#1}}
%\providecommand{\partitle}[1]{\textcolor{alertcolor}{#1}}
%\providecommand{\light}[1]{#1}
%
%\usepackage{mleftright}
%\usepackage{dsfont}
%
%\providecommand{\vrtr}{{\vr{\Vtheta_r}}}
%\providecommand{\pdvr}[2]{\frac{\partial }{\partial #2}#1}
%
%\providecommand\drawNodes[3]{
%  % #1 (str): namespace
%  % #2 (list[list[str]]): list of labels to print in the node of each neuron
%  \foreach\neurons [count=\lyrIdx] in #2 {
%    \StrCount{\neurons}{,}[\lyrLength] % use xstring package to save each layer size into \lyrLength macro
%    \foreach\n [count=\nIdx] in \neurons {
%      \node[neuron={#3}] (#1-\lyrIdx-\nIdx) at (2*\lyrIdx, \lyrLength/2-1.4*\nIdx) {\n};
%    }
%  }
%}
%
%\providecommand{\fq}{\func{q}}
%\providecommand{\fqi}{\func{q_i}}
%\providecommand{\fqj}{\func{q}}
%\providecommand{\fqu}{\func{q_1}}
%\providecommand{\fqn}{\func{q_n}}
%\providecommand{\fqmj}{\func{q_{-j}}}
%\providecommand{\fQ}{\func{\mathcal{Q}}}
%\providecommand{\um}[1]{\scalebox{0.4}[1.0]{\( - \)}#1}
%\providecommand{\ELBO}{\mathrm{ELBO}}
%
%\providecommand{\va}[1]{\text{\textit{Var}}[#1]}
%
%% ---------- Macros ----------
%
%\providecommand{\calL}{\mathcal{L}}
%\providecommand{\calR}{\mathcal{R}}
%\providecommand{\calT}{\mathcal{T}}
%\providecommand{\calY}{\mathcal{Y}}
%\providecommand{\calX}{\mathcal{X}}
%\providecommand{\calH}{\mathcal{H}}
%\providecommand{\EmpR}{\overline{\mathcal{R}}}
%
%\definecolor{torblue}{RGB}{31,60,110}
%\definecolor{toraccent}{RGB}{0,112,192}
%\definecolor{torgray}{RGB}{220,225,235}
%\definecolor{alertred}{RGB}{180,30,30}
%\definecolor{exgreen}{RGB}{0,128,64}
%
%\definecolor{cerise}{HTML}{DA3A62}
%\definecolor{teal}{HTML}{029386}
%\definecolor{slateblue}{HTML}{5B5EA6}
%\definecolor{orange}{HTML}{F97306}

% ==============================================================================

%\leadauthor{Giorgio Gambosi}
%\shorttitle{Notes on Machine Learning}

\usepackage{amsmath, amssymb, mathtools}
\usepackage{booktabs}
\usepackage{tikz}
\usetikzlibrary{arrows.meta, positioning, shapes.geometric}
\usepackage{xcolor}
\usepackage{microtype}
\usepackage{graphicx}
% indicator function via amsmath

% ---------- Macros ----------

\providecommand{\calX}{\mathcal{X}}
\providecommand{\calY}{\mathcal{Y}}
\providecommand{\calH}{\mathcal{H}}
\providecommand{\calT}{\mathcal{T}}
\providecommand{\calA}{\mathcal{A}}
\providecommand{\calR}{\mathcal{R}}
\providecommand{\barR}{\overline{\mathcal{R}}}
\providecommand{\pM}{p_M}
\providecommand{\pC}{p_C}
\providecommand{\indicator}[1]{\mathbf{1}\!\left[#1\right]}
\providecommand{\argmax}{\operatorname*{arg\,max}}
\providecommand{\argmin}{\operatorname*{arg\,min}}

\definecolor{cerise}{HTML}{DA3A62}
\definecolor{teal}{HTML}{029386}
\definecolor{slateblue}{HTML}{5B5EA6}
\definecolor{orange}{HTML}{F97306}

% ---------- Macros ----------

\providecommand{\calL}{\mathcal{L}}
\providecommand{\calR}{\mathcal{R}}
\providecommand{\calT}{\mathcal{T}}
\providecommand{\calY}{\mathcal{Y}}
\providecommand{\calX}{\mathcal{X}}
\providecommand{\calH}{\mathcal{H}}
\providecommand{\EmpR}{\overline{\mathcal{R}}}

\usepackage{amsmath,amssymb,mathtools}
\usepackage{booktabs}
\usepackage{graphicx}
\usepackage{xcolor}
\usepackage{tikz}
\usetikzlibrary{arrows.meta,positioning}
\usepackage{listings}
\usepackage{array}
\usepackage{multirow}

\lstset{
  language=Python,
  basicstyle=\ttfamily\tiny,
  keywordstyle=\color{MainBlue}\bfseries,
  commentstyle=\color{gray},
  frame=single,
  breaklines=true,
}

% ---------- Macros ----------
\providecommand{\bth}{\boldsymbol{\theta}}
\providecommand{\bx}{\mathbf{x}}
\providecommand{\bw}{\mathbf{w}}
\providecommand{\calT}{\mathcal{T}}
\providecommand{\calL}{\mathcal{L}}
\providecommand{\Rbar}{\overline{\mathcal{R}}}
\providecommand{\defeq}{\stackrel{\text{def}}{=}}
\usepackage{amsmath,amssymb,bm}
\usepackage{pgfplots}
\pgfplotsset{compat=1.18}
\usepackage{tikz}
\usetikzlibrary{arrows.meta,positioning,shapes,decorations.pathreplacing,calc,decorations.markings}
\usepgfplotslibrary{fillbetween}
\usepackage{booktabs}
\usepackage{multicol}
\usepackage{graphicx}
\usepackage{xcolor}
\usepackage{mathtools}

% ── Math macros ──────────────────────────────────────────────────────────────


\providecommand{\ow}{\overline{\mathbf{w}}}
\providecommand{\ox}{\overline{\mathbf{x}}}
\providecommand{\oX}{\overline{\mathbf{X}}}

\usepackage{amsmath, amssymb, mathtools}
\usepackage{booktabs}
\usepackage{tikz}
\usetikzlibrary{arrows.meta, positioning, shapes.geometric}
\usepackage{xcolor}
\usepackage{microtype}
\usepackage{graphicx}
% indicator function via amsmath

% ---------- Macros ----------

\providecommand{\calX}{\mathcal{X}}
\providecommand{\calY}{\mathcal{Y}}
\providecommand{\calH}{\mathcal{H}}
\providecommand{\calT}{\mathcal{T}}
\providecommand{\calA}{\mathcal{A}}
\providecommand{\calR}{\mathcal{R}}
\providecommand{\barR}{\overline{\mathcal{R}}}
\providecommand{\pM}{p_M}
\providecommand{\pC}{p_C}
\providecommand{\indicator}[1]{\mathbf{1}\!\left[#1\right]}
\providecommand{\argmax}{\operatorname*{arg\,max}}
\providecommand{\argmin}{\operatorname*{arg\,min}}

\usepackage{amsmath,amsfonts,amssymb,mathtools}
\usepackage{bm}
\usepackage{upgreek}

% ---- Graphics -------------------------------------------
\usepackage{graphicx}
\usepackage{microtype}
% ---- tcolorbox ------------------------------------------
\usepackage[skins]{tcolorbox}

\providecommand{\Ii}{\mathbf{I}}
\providecommand{\bZero}{\mathbf{0}}
\providecommand{\bX}{\mathbf{X}}
\providecommand{\bZ}{\mathbf{Z}}
\providecommand{\btt}{\mathbf{t}}
\providecommand{\bS}{\mathbf{S}}
\providecommand{\bx}{\mathbf{x}}
\providecommand{\bz}{\mathbf{z}}
\providecommand{\bw}{\mathbf{w}}
\providecommand{\ox}{\overline{\mathbf{x}}}
\providecommand{\ow}{\overline{\mathbf{w}}}
\providecommand{\oX}{\overline{\mathbf{X}}}
\providecommand{\bPsi}{\boldsymbol{\Psi}}
\providecommand{\halert}[1]{{\color{AlertRed}\textbf{#1}}}

\providecommand{\SW}{\mathbf{S}_W}
\providecommand{\SB}{\mathbf{S}_B}
\providecommand{\ST}{\mathbf{S}_T}
\providecommand{\Smat}{\mathbf{S}}
\providecommand{\GX}{G_{\mathbf{X}}}

\definecolor{torred}{RGB}{153,0,0}
\definecolor{torcyan}{RGB}{0,112,192}

\providecommand{\ow}{\overline{\mathbf{w}}}
\providecommand{\ox}{\overline{\mathbf{x}}}
\providecommand{\oX}{\overline{\mathbf{X}}}

\providecommand{\Hmat}{\mathbf{H}}

\providecommand{\mlalert}[1]{{\color{AccentBlue}\textbf{#1}}}
\providecommand{\mlemph}[1]{{\color{MainBlue}\textit{#1}}}

\setbeamertemplate{theorems}[numbered]
%\newtheorem{mytheorem}{Theorem}
%\newtheorem{mydefinition}{Definition}

% ---------- Macros ----------

\providecommand{\calX}{\mathcal{X}}
\providecommand{\calY}{\mathcal{Y}}
\providecommand{\calH}{\mathcal{H}}
\providecommand{\calT}{\mathcal{T}}
\providecommand{\calA}{\mathcal{A}}
\providecommand{\calR}{\mathcal{R}}
\providecommand{\calF}{\mathcal{F}}
\providecommand{\barR}{\overline{\mathcal{R}}}
\providecommand{\pM}{p_M}
\providecommand{\pC}{p_C}
\providecommand{\indicator}[1]{\mathbf{1}\!\left[#1\right]}
\providecommand{\vcd}[1]{\operatorname{VCdim}(#1)}
\providecommand{\Prob}[2]{\mathbb{P}_{#1}\!\left[#2\right]}
% ---------- Theme ----------
% ---------- Custom colours (matching PDF model) ----------
\definecolor{TvDark}{RGB}{30,50,80}        % dark navy for header bar
\definecolor{TvMid}{RGB}{70,100,140}       % medium blue for blocks
\definecolor{TvLight}{RGB}{210,220,235}    % pale blue for backgrounds
\definecolor{TvAlert}{RGB}{160,30,30}      % dark red for alerts / warnings
\definecolor{TvGreen}{RGB}{20,100,50}      % green for examples

\providecommand{\SW}{\mathbf{S}_W}
\providecommand{\SB}{\mathbf{S}_B}
\providecommand{\ST}{\mathbf{S}_T}
\providecommand{\Smat}{\mathbf{S}}
\providecommand{\GX}{G_{\mathbf{X}}}

\usepackage{subcaption}

\newcommand{\blankpage}{
\clearpage
\thispagestyle{empty}
\null
\clearpage
}